现在证明 Theorem~\ref{thm:global-boundary-regularity}.

\begingroup
\renewcommand{\endProof}{\par\nobreak\hfill$\blacksquare$\par\vspace{0.1cm}}
\begin{Proof}
现在只需将各部分结合起来, 即可完成整体正则性定理的证明. 对每个点 $\overline{\bm{x}}\in\partial\Omega$, 使用 Lemma~\ref{lem:local-boundary-regularity-transfer} 中构造的邻域 $\widetilde Q$, 再添入一个集合 $\Omega_0\subset\subset\Omega$ 覆盖内部, 就给出了 $\overline\Omega$ 的开覆盖. 由于 $\overline\Omega$ 紧, 存在有限子覆盖 (假设其中包含 $\Omega_0$, 并记为 $\{\Omega_i\}_{i=0}^N$), 使得
\begin{equation}\label{eq:global-boundary-regularity-finite-open-cover}
\Omega\subset\bigcup_{i=0}^N\Omega_i.
\end{equation}
现在对 $\Omega_0$ 利用内部正则性结果 (Theorem~\ref{thm:second-order-interior-regularity}), 并对其他每个集合利用 Lemma~\ref{lem:local-boundary-regularity-transfer}, 得到
\begin{equation}\label{eq:global-boundary-regularity-summed-local-estimates}
\|u\|_{H^2(\Omega)}\leq\sum_{i=0}^N\|u\|_{H^2(\Omega_i)}\leq C(\|f\|_{L^2(\Omega)}+\|u\|_{H^1(\Omega)}).
\end{equation}
标准地应用 Ehrling 引理 (Theorem~\ref{thm:ehrling-compact-intermediate-norm}), 即得最终结果.
\end{Proof}
\endgroup
